\chapter{Environnement et construction}\label{chap:environnement}

\section{Versions retenues}
Le tableau~\ref{tab:versions} liste les composants et la raison de chaque version. Le sujet fournit une liste de
dépendances ; plusieurs entrées ont dû être corrigées pour fonctionner avec le JDK et Tomcat installés.

\begin{table}[htbp]
\centering
\caption{Composants et versions}\label{tab:versions}
\small
\begin{tabularx}{\linewidth}{@{}llX@{}}
\toprule
\textbf{Composant} & \textbf{Version} & \textbf{Remarque}\\
\midrule
JDK (exécution) & 17 & Tomcat et Eclipse l'utilisent.\\
Compilation & \code{release} 11 & EclipseLink 2.7.5 ne lit pas les classes Java 17 (section~\ref{sec:asm}).\\
Tomcat & 9.0.123, port 8081 & Espace de noms \code{javax.*}, compatible JSF 2.2.\\
JSF (Mojarra) & 2.2.20 & Tomcat ne fournit pas JSF : \code{jsf-api} et \code{jsf-impl} sont embarqués dans le WAR.\\
CDI (Weld) & 3.1.9.Final & Version \textit{servlet shaded} pour Tomcat.\\
JPA (EclipseLink) & 2.7.5 & Apporte l'API \code{javax.persistence} 2.2.\\
Pilote MySQL & 8.0.33 & Remplace le 8.0.12 du sujet, compatible MySQL 8.\\
Servlet API & 4.0.1, \code{provided} & Fournie par Tomcat ; ne doit pas se trouver dans le WAR.\\
JAXB API & 2.3.1 & Absente du JDK depuis Java 11 ; requise par Mojarra (chapitre~\ref{chap:incidents}).\\
\bottomrule
\end{tabularx}
\end{table}

\section{Corrections apportées aux dépendances du sujet}
\begin{itemize}
  \item \textbf{Servlet API} : le sujet donne la version 4.0.0 sans portée. Sans \code{provided}, l'API se retrouve
        dans le WAR et entre en conflit avec celle de Tomcat.
  \item \textbf{JPA} : l'artefact \code{org.eclipse.persistence:javax.persistence:2.0.0} du sujet n'est \emph{pas}
        déclaré. EclipseLink 2.7.5 apporte déjà l'API 2.2 ; deux versions de l'API sur le même chemin de classes
        provoqueraient des conflits.
  \item \textbf{JSTL et JSP} : Facelets n'en a pas besoin (les balises \code{c:} et \code{fn:} sont implémentées dans
        Mojarra). Les déclarations de l'atelier~2 sont donc retirées.
  \item \textbf{JAXB} : ajouté après l'incident du chapitre~\ref{chap:incidents}.
\end{itemize}

\lstinputlisting[language=XML,caption={Extrait du \texttt{pom.xml} : JSF et JAXB},label={lst:pom}]{extraits/pom-jsf-jaxb.txt}

\section{Configuration de JSF}
Le contrôleur frontal est déclaré dans \code{web.xml} et associé au motif \code{*.xhtml}. Ce choix garde des adresses
lisibles (\code{vitrine.xhtml?categorie=2}) et signifie qu'une page Facelets est atteinte directement par son
nom, sans préfixe \code{/faces/}.

\lstinputlisting[language=XML,caption={Déclaration de \texttt{FacesServlet}},label={lst:faces}]{extraits/web-facesservlet.txt}

Le même fichier configure la phase de développement (\code{PROJECT\_STAGE=Development}, messages d'erreur plus
détaillés), la sauvegarde de l'état de la vue sur le serveur, l'ordre des filtres (sécurité puis authentification),
les pages d'erreur (403, 404, 405, 500 et toute exception) et un bloc de sécurité interdisant l'accès direct à
\code{/resources/components/*}.

\section{Construction et démarrage}
La construction se fait par Maven (\code{mvn -q clean package}), le déploiement par Eclipse sur Tomcat. Au démarrage,
la console doit montrer, dans l'ordre : l'initialisation de Weld, celle de Mojarra, puis la connexion EclipseLink
(\code{login successful}). La figure~\ref{fig:demarrage} en donne la trace.

\fig{05-tomcat-demarrage}{Démarrage de Tomcat : Weld, Mojarra et connexion EclipseLink}{fig:demarrage}

\section{Compiler pour Java 11 sous Java 17}\label{sec:asm}
EclipseLink 2.7.5 lit les classes des entités avec une version d'ASM qui rejette les fichiers de classe plus récents
que Java 13 (\guillemotleft~Unsupported class file major version 61~\guillemotright{}). Compilées pour Java 17, les entités ne seraient
pas découvertes et l'application démarrerait sans table. La propriété \code{maven.compiler.release=11} compile
pour Java 11 tout en gardant le JDK 17 : les classes ainsi produites s'exécutent normalement sous Java 17. Cette
contrainte était déjà apparue à l'atelier~2 ; elle est ici reconduite et documentée dans le \code{pom.xml}.
